%=======================================================================
%  CSE558 Data Science - Midterm Project Report Template
%  Monsoon 2026, IIIT-Delhi
%
%  How to use this template
%   1. Set your project title in the EDIT THIS block below.
%   2. Replace everything in LIGHT BLUE with your own text and numbers.
%      The numbers in the example are made up for illustration.
%   3. Grey italic lines are instructions. Delete them as you write, or
%      hide them all at once with \showhintsfalse.
%   4. Upload your plots to the project and put their file names in
%      \figfile{...}{...}. Until a file exists, a labelled box is shown.
%
%  Sections 1-5 follow items 1-5 of the report guidelines; Section 6 is
%  the limitations of your analysis. The link to your code (item 6) goes
%  at the top. At most 8 pages.
%  Compile with pdfLaTeX (Overleaf's default setting works), twice.
%=======================================================================
\documentclass[10pt,a4paper,twocolumn]{article}

%========================= EDIT THIS =========================
% Your project title: it sets the heading AND the PDF's title.
\newcommand{\projecttitle}{Do crowded neighbourhoods drag ratings down?
  Evidence from Bengaluru restaurant listings}
\newif\ifshowhints
\showhintstrue      % \showhintsfalse hides every grey instruction line
%=============================================================

%-------------------------------------------- fonts (no need to edit)
\usepackage{iftex}
\ifPDFTeX
  \usepackage[T1]{fontenc}
  \usepackage[utf8]{inputenc}
  \IfFileExists{carlito.sty}%
    {\usepackage[sfdefault]{carlito}}%              Calibri look
    {\usepackage[scaled=0.92]{helvet}\renewcommand{\familydefault}{\sfdefault}}%
  \usepackage{courier}
\else
  \usepackage[no-math]{fontspec}
  \setmainfont{Carlito}[Extension=.ttf, UprightFont=*-Regular, BoldFont=*-Bold,
    ItalicFont=*-Italic, BoldItalicFont=*-BoldItalic]
\fi

%----------------------------------------------------------- packages
\usepackage[margin=18mm,bottom=20mm,columnsep=7mm]{geometry}
\usepackage{amsmath}
\usepackage{xcolor}
\usepackage{graphicx}
\usepackage{tabularx}
\usepackage{array}
\usepackage{booktabs}
\usepackage{enumitem}
\usepackage{titlesec}
\usepackage{needspace}
\usepackage[font=small,labelfont=bf]{caption}
\usepackage[hidelinks]{hyperref}
\hypersetup{pdftitle={\projecttitle},
            pdfsubject={CSE558 Data Science, Monsoon 2026: Midterm Project Report}}

%------------------------------------------------------------- layout
\definecolor{eg}{HTML}{7396C4}         % worked example (light blue)
\definecolor{hintgrey}{gray}{0.58}     % instructions (light grey)

\titleformat{\section}{\normalfont\large\bfseries}{\thesection.}{0.5em}{}
\titlespacing*{\section}{0pt}{12pt}{3pt}
\titleformat{\subsection}{\normalfont\normalsize\bfseries}{\thesubsection}{0.5em}{}
\titlespacing*{\subsection}{0pt}{10pt}{2pt}
\newcommand{\sectionbreak}{\needspace{10\baselineskip}}
\newcommand{\subsectionbreak}{\needspace{10\baselineskip}}

\setlength{\parindent}{0pt}
\setlength{\parskip}{4pt}
\setlist{leftmargin=1.2em, itemsep=2pt, topsep=2pt, parsep=0pt}
\renewcommand{\arraystretch}{1.25}
\setlength{\tabcolsep}{4pt}
\newcolumntype{L}[1]{>{\raggedright\arraybackslash}p{#1}}
\newcolumntype{Y}{>{\raggedright\arraybackslash}X}
\newcolumntype{K}[1]{>{\raggedright\arraybackslash\bfseries\boldmath}p{#1}}  % bold labels

%--------------------------------------------------------- shorthands
% Instruction line. Delete these as you fill the template in.
\newcommand{\hint}[1]{\ifshowhints{\footnotesize\itshape\color{hintgrey}#1\par}\fi}
% Worked example: inline or in a table cell. REPLACE the contents.
\newcommand{\eg}[1]{\textcolor{eg}{#1}}
% Worked example running over a whole paragraph.
\newcommand{\egblock}[1]{{\color{eg}#1\par}}
% Bold label followed by an example.
\newcommand{\field}[2]{\textbf{#1}~\eg{#2}\par}
% Part a)-c) of a test, and a small bold heading.
\newcommand{\tpart}[2]{\par\needspace{12\baselineskip}\smallskip\textbf{#1)~#2}\par}
\newcommand{\minihead}[1]{\par\needspace{8\baselineskip}\smallskip\textbf{#1}\par}
% Column and file names.
\newcommand{\code}[1]{\texttt{#1}}
% \figfile{file}{height}: your plot, or a labelled box until the file exists.
\newcommand{\figfile}[2]{%
  \IfFileExists{#1}%
    {\includegraphics[width=\linewidth,height=#2,keepaspectratio]{#1}}%
    {\fbox{\parbox[c][#2][c]{\dimexpr\linewidth-2\fboxsep-2\fboxrule\relax}%
      {\centering\small\itshape\color{hintgrey}Your plot: \detokenize{#1}}}}}

%=======================================================================
\begin{document}

\twocolumn[%
{\small\bfseries CSE558 Data Science\ \ \textbullet\ \ Midterm Project Report\ \ \textbullet\ \ Monsoon 2026\par}
\vspace{4pt}
{\LARGE\bfseries\raggedright\projecttitle\par}
\vspace{6pt}
\rule{\linewidth}{1pt}

\hint{Everything in light blue is a worked example: replace it with your own. Its numbers are
made up for illustration. Grey italic lines are instructions: delete them, or hide them all with
\texttt{\textbackslash showhintsfalse} at the top of the file, where you also set your project
title. Sections 1--5 follow items 1--5 of the report guidelines; Section 6 is the limitations of
your analysis. At most 8 pages.}

{\small
\begin{tabularx}{\linewidth}{@{}K{2.4cm} Y@{}}
\toprule
Group name & \eg{Team Kurtosis} \\
Dataset & \eg{kaggle.com/datasets/himanshupoddar/zomato-bangalore-restaurants} \\
Code & \eg{github.com/team-kurtosis/cse558-midterm (notebook and README)} \\
\bottomrule
\end{tabularx}}

\hint{Code: a link to the code written by you and a readme file, which anyone with the link can
open. Also upload your notebook or script on the classroom.}

\vspace{2pt}
{\small
\begin{tabularx}{\linewidth}{@{}L{0.5cm} L{3.4cm} L{2.3cm} L{4.8cm} Y@{}}
\toprule
\textbf{\#} & \textbf{Name} & \textbf{Roll no.} & \textbf{Institute email} & \textbf{Contribution} \\
\midrule
1 & \eg{Aditi Rao} & \eg{2022145} & \eg{aditi22145@iiitd.ac.in} & \eg{Dataset and size check (1)} \\
2 & \eg{Rohan Mehta} & \eg{2022311} & \eg{rohan22311@iiitd.ac.in} & \eg{Preparation, before and after (2, 3)} \\
3 & \eg{Sneha Iyer} & \eg{2022078} & \eg{sneha22078@iiitd.ac.in} & \eg{EDA and problem statements (4)} \\
4 & \eg{Karan Bose} & \eg{2022260} & \eg{karan22260@iiitd.ac.in} & \eg{Test 1 (5)} \\
5 & \eg{Meera Nair} & \eg{2022019} & \eg{meera22019@iiitd.ac.in} & \eg{Test 2, limitations, code and README (5, 6)} \\
\bottomrule
\end{tabularx}}
\vspace{12pt}
]

%-----------------------------------------------------------------------
\section{Dataset}
\hint{Where you have collected the dataset, how many rows and columns there are, and whether
\#instances $\times$ \#attributes is at least 3M after encoding (show the arithmetic). Name some
features with their interpretation to the real world.}

\field{Source:}{Kaggle, ``Zomato Bangalore Restaurants'': scraped restaurant listings,
downloaded on 28 August 2026. One row is one outlet.}

\field{What is raw about it:}{ratings are stored as text (\code{"4.1/5"}, \code{"NEW"},
\code{"-"}); the cost for two is a string with thousands separators; \code{dish\_liked} is a
comma-separated list; some chain outlets are listed twice; 15\% of ratings are missing or
\code{NEW}.}

{\small
\begin{tabularx}{\linewidth}{@{}L{3.0cm} Y Y@{}}
\toprule
 & \textbf{Raw file} & \textbf{After encoding} \\
\midrule
Instances (rows) & \eg{51{,}717} & \eg{48{,}304} \\
Attributes (columns) & \eg{17} & \eg{96} \\
Instances $\times$ attributes & \eg{879{,}189} & \eg{4{,}637{,}184} \\
\bottomrule
\end{tabularx}}

\field{Size check:}{$48{,}304\times96=4{,}637{,}184\ge3{,}000{,}000$, so the requirement is met.
One-hot encoding of \code{location} and \code{cuisines} is what takes it past 3M.}

\minihead{Key features}
\hint{Some key features, with their type (nominal, discrete or continuous) and real-world
meaning. Include the response variable, if there is one.}
{\footnotesize
\begin{tabularx}{\linewidth}{@{}L{2.2cm} L{1.5cm} Y@{}}
\toprule
\textbf{Feature} & \textbf{Type} & \textbf{Real-world meaning} \\
\midrule
\eg{\code{rate}} & \eg{continuous} & \eg{mean customer rating out of 5; our response} \\
\eg{\code{votes}} & \eg{discrete} & \eg{number of reviews behind the rating} \\
\eg{\code{location}} & \eg{nominal} & \eg{Bengaluru neighbourhood, 93 levels} \\
\eg{\code{cuisines}} & \eg{nominal} & \eg{cuisine tags, several per outlet} \\
\eg{\code{cost\_two}} & \eg{continuous} & \eg{typical bill for two, in rupees} \\
\eg{\code{online\_order}} & \eg{nominal} & \eg{takes online orders (yes/no)} \\
\eg{\code{book\_table}} & \eg{nominal} & \eg{takes table bookings (yes/no)} \\
\bottomrule
\end{tabularx}}

%-----------------------------------------------------------------------
\section{Data preparation}
\hint{How did you prepare your data (encoding, outliers, missing values, duplicates), and why do
you follow such steps? One row per step; the ``why'' matters most.}

{\footnotesize
\begin{tabularx}{\linewidth}{@{}L{1.9cm} L{2.3cm} Y@{}}
\toprule
\textbf{Issue found} & \textbf{What we did} & \textbf{Why, and not the alternative} \\
\midrule
\eg{\code{rate} is text, with \code{NEW} and \code{-}}
 & \eg{parsed to a number; 3{,}359 \code{NEW} rows dropped; \code{-} set to missing}
 & \eg{\code{NEW} means not yet rated, not rated zero. Coding it as 0 would invent a group of
   very bad restaurants.} \\
\eg{4{,}399 ratings still missing}
 & \eg{filled with the median rating of the same location, and flagged as imputed}
 & \eg{keeps the rows for later modelling; the flag lets our tests in Section 5 leave the
   filled-in values out} \\
\eg{54 duplicate chain rows}
 & \eg{dropped on (name, address)}
 & \eg{an outlet listed twice would inflate exactly the crowded areas we study} \\
\eg{\code{cost\_two} stored as ``1,200''}
 & \eg{separators removed, converted to a number}
 & \eg{a format fix; nothing is lost} \\
\eg{612 costs above the IQR fence}
 & \eg{capped at the fence ($Q_3+1.5\,\mathrm{IQR}$)}
 & \eg{they are real luxury outlets, so we keep the rows but stop them dominating means} \\
\eg{\code{location} has 93 levels}
 & \eg{one-hot encoded; 5 levels with under 30 rows merged into \code{Other}}
 & \eg{rare levels give unstable estimates; merging keeps every crowded area intact} \\
\bottomrule
\end{tabularx}}

%-----------------------------------------------------------------------
\section{What changed after preparation}
\hint{What are the changes you got after data preparation? Give before-and-after observations.}

{\footnotesize
\begin{tabularx}{\linewidth}{@{}L{2.9cm} L{1.5cm} Y@{}}
\toprule
\textbf{Measure} & \textbf{Before} & \textbf{After} \\
\midrule
Rows & \eg{51{,}717} & \eg{48{,}304} \\
Attributes & \eg{17} & \eg{96} \\
Missing or unusable ratings & \eg{7{,}758 (15.0\%)} & \eg{0: 3{,}359 dropped, 4{,}399 imputed and flagged} \\
Duplicate rows & \eg{54} & \eg{0} \\
\code{cost\_two} above the IQR fence & \eg{612} & \eg{0 (capped at the fence)} \\
\bottomrule
\end{tabularx}}

\egblock{Parsing \code{rate} turned the response from text into a number, which is what makes
any test possible. Encoding \code{location} and \code{cuisines} took the attributes from 17 to
96, past the size requirement. We lost 3{,}413 rows: 3{,}359 unrated new listings and 54
duplicates.}

%-----------------------------------------------------------------------
\section{Basic questions and problem statement}
\hint{What are your basic questions about your data, and how is your problem statement novel
from them? Show a few EDA plots, each with a one-line takeaway.}

\begin{figure}[!ht]
  \figfile{fig_ratings.pdf}{3.6cm}
  \caption{\eg{Ratings bunch between 3.5 and 4.2 and trail off to the left (skewness $-0.7$);
  31\% of outlets reach 4.0.}}
\end{figure}

\begin{figure}[!ht]
  \figfile{fig_density.pdf}{3.6cm}
  \caption{\eg{Outlets per neighbourhood run from 32 to 3{,}660, so a few areas would dominate
  a straight-line fit on density.}}
\end{figure}

\minihead{Basic questions}
\begin{enumerate}[label=\arabic*.]
  \item \eg{Do costlier restaurants get higher ratings?}
  \item \eg{Which cuisine has the best average rating?}
  \item \eg{Do outlets with more votes have higher ratings?}
\end{enumerate}

\minihead{Our problem statements}
\hint{One or two sentences each, specific to your features, and phrased so each could come out
either way. Mark one as primary.}

\noindent\fbox{\begin{minipage}{\dimexpr\linewidth-2\fboxsep-2\fboxrule\relax}
\textbf{P1 (primary).} \eg{Among restaurants with comparable cost and cuisine, does rating fall
as the number of competing outlets in the same neighbourhood rises?}

\smallskip
\textbf{P2.} \eg{Within the most crowded neighbourhoods, do outlets that take online orders rate
higher than those that do not?}
\end{minipage}}

\field{How they are novel:}{the basic questions look at each outlet on its own. Ours depend on
the neighbourhood an outlet competes in, which no single row records, so they need a derived
feature (outlets per neighbourhood) and a like-for-like comparison.}

%-----------------------------------------------------------------------
\section{Statistical inferences}
\hint{Conduct 2--3 hypothesis tests on your problem statement(s), each with parts a)--c) below.
Know your tests and the reason for using them, and compare them with other tests.}

\field{Rows tested:}{the 43{,}905 outlets with an observed rating; the 4{,}399 imputed ratings
are left out. Comparable outlets, North Indian at INR 400--800 for two: 7{,}877. We rank them by
the number of outlets in their neighbourhood and cut them into thirds at neighbourhood
boundaries: the most crowded third $D$ (2{,}805 outlets in 9 neighbourhoods), a middle third
(not used) and the least crowded third $S$ (2{,}496 outlets in 51 neighbourhoods).}

%------------------------------------------------------------ Test 1
\subsection*{Test 1 (P1): do comparable outlets rate lower in the most crowded neighbourhoods?}

\tpart{a}{Hypotheses, written before running the test}
\hint{Before running the test, write the null and alternative hypotheses ($H_0$ carries the
equality), the significance level, one- or two-sided, the test you will use, and the smallest
difference that would matter in the real world.}
{\small
\begin{tabularx}{\linewidth}{@{}K{3.4cm} Y@{}}
\toprule
$H_0$ & \eg{$\mu_D-\mu_S\ge0$} \\
$H_1$ & \eg{$\mu_D-\mu_S<0$} \\
Significance level & \eg{$\alpha=0.05$} \\
Side & \eg{one-sided (left)} \\
Test & \eg{Welch t-test} \\
Smallest difference that matters & \eg{0.10 star} \\
\bottomrule
\end{tabularx}}

\egblock{$\mu_D$, $\mu_S$: mean rating of comparable outlets in the most and least crowded
thirds. One-sided, because P1 claims that ratings fall. 0.10 star, because Zomato shows ratings
to one decimal: a smaller gap is invisible to a diner.}

\tpart{b}{Other tests we considered}
\hint{Which other tests did you consider, and why did you not use them? Say first why your test
fits your data.}
\egblock{Two independent groups, a continuous response and an unknown $\sigma$ call for a
two-sample t-test. We use Welch's version because it does not assume equal spreads.}
{\footnotesize
\begin{tabularx}{\linewidth}{@{}L{2.4cm} Y@{}}
\toprule
\textbf{Also considered} & \textbf{Why we did not use it} \\
\midrule
\eg{Two-sample z-test} & \eg{needs the population $\sigma$; we only have sample spreads} \\
\eg{Pooled t-test} & \eg{assumes equal variances, which an F-test rejects ($F=1.31$, $p<0.001$). Run anyway, it gives $t=-4.96$ and the same decision.} \\
\eg{Paired t-test} & \eg{nothing to pair: each outlet is in one group} \\
\eg{t-test on the correlation $r$} & \eg{sees only a straight-line trend, which Figure 2 shows a few large areas would drive} \\
\bottomrule
\end{tabularx}}

\tpart{c}{Result}
\hint{Report the test statistic, degrees of freedom, p-value and decision, and the estimate with
a 95\% confidence interval. Is the difference large enough to matter?}
{\small
\begin{tabularx}{\linewidth}{@{}K{3.4cm} Y@{}}
\toprule
Estimate & \eg{$-0.06$ star (3.68 in $D$, 3.74 in $S$)} \\
95\% confidence interval & \eg{$-0.084$ to $-0.036$} \\
Test statistic (d.f.) & \eg{$t=-4.92$ (4{,}984)} \\
p-value & \eg{$4.4\times10^{-7}$} \\
Decision & \eg{reject $H_0$} \\
\bottomrule
\end{tabularx}}

\egblock{Comparable outlets in the most crowded neighbourhoods rate 0.06 star lower on average.
The difference is real, but it does not matter by our standard: the whole interval lies inside
$\pm0.10$ star.}

%------------------------------------------------------------ Test 2
\subsection*{Test 2 (P2): does online ordering lift ratings in the most crowded neighbourhoods?}
\hint{The same three parts as Test 1.}

\tpart{a}{Hypotheses, written before running the test}
{\small
\begin{tabularx}{\linewidth}{@{}K{3.4cm} Y@{}}
\toprule
$H_0$ & \eg{$\mu_{\text{on}}-\mu_{\text{off}}\le0$} \\
$H_1$ & \eg{$\mu_{\text{on}}-\mu_{\text{off}}>0$} \\
Significance level & \eg{$\alpha=0.05$} \\
Side & \eg{one-sided (right)} \\
Test & \eg{Welch t-test} \\
Smallest difference that matters & \eg{0.10 star} \\
\bottomrule
\end{tabularx}}

\egblock{$\mu_{\text{on}}$, $\mu_{\text{off}}$: mean rating, within $D$, of outlets that do and
do not take online orders. One-sided, because P2 claims a lift.}

\tpart{b}{Other tests we considered}
\egblock{Again two independent groups, a continuous response and an unknown $\sigma$, so a
two-sample t-test; we fixed Welch's version in advance.}
{\footnotesize
\begin{tabularx}{\linewidth}{@{}L{2.4cm} Y@{}}
\toprule
\textbf{Also considered} & \textbf{Why we did not use it} \\
\midrule
\eg{Two-sample z-test} & \eg{needs the population $\sigma$; we only have sample spreads} \\
\eg{Pooled t-test} & \eg{also valid here, since an F-test does not reject equal variances ($F=1.10$, $p=0.08$); it gives $t=1.25$ and the same decision. We keep Welch, fixed in advance.} \\
\eg{Paired t-test} & \eg{no outlet is in both groups} \\
\bottomrule
\end{tabularx}}

\tpart{c}{Result}
{\small
\begin{tabularx}{\linewidth}{@{}K{3.4cm} Y@{}}
\toprule
Estimate & \eg{$+0.02$ star (3.69 with online orders, 3.67 without)} \\
95\% confidence interval & \eg{$-0.012$ to $+0.052$} \\
Test statistic (d.f.) & \eg{$t=1.24$ (2{,}086)} \\
p-value & \eg{0.11} \\
Decision & \eg{fail to reject $H_0$} \\
\bottomrule
\end{tabularx}}

\egblock{Our data gives no evidence that taking online orders lifts ratings in crowded
neighbourhoods. The interval rules out a lift larger than 0.052 star, so any lift is below the
0.10 star that would matter.}

\hint{Running a third test? Add it here as Test 3, with the same three parts: copy the Test 2
block.}

%-----------------------------------------------------------------------
\section{Limitations}
\hint{What your analysis cannot show, and why: rows you dropped, assumptions that hold only
approximately, groups your conclusions do not cover, and association versus cause.}

\begin{itemize}
  \item \eg{\textbf{Association, not cause.} Crowded neighbourhoods also differ in rent, footfall
  and customers; our tests cannot separate density from these.}
  \item \eg{\textbf{Outlets are not fully independent.} Outlets in one neighbourhood share
  conditions, so our effective sample is smaller than 5{,}301 and our p-values are optimistic.
  A test on the 60 neighbourhood averages would be stricter.}
  \item \eg{\textbf{Established outlets only.} The 3{,}359 unrated new listings were dropped in
  Section 2.}
  \item \eg{\textbf{Imputed ratings left out.} The 4{,}399 outlets with a missing rating may
  differ from the rest, and they are not in our tests.}
  \item \eg{\textbf{One cuisine and price band.} We compared North Indian outlets at INR
  400--800 for two; other cuisines and prices may behave differently.}
\end{itemize}

\end{document}
